\subsection*{一步一步建立直观}

\paragraph{1. 把抽象条件换成可追踪的动作}

面对线性方程，先把未知数看作输入向量，把方程右侧看作期望输出。右图让矩阵两列成为基向量的去向，于是求解问题变成“哪一个输入能到达目标”。选择合适表示常常比立即展开计算更重要，因为它能先暴露问题的结构。



\paragraph{2. 检查退化情形，寻找缺失的条件}

把两列调到共线，输出平面压成直线。此时某些目标不可达，某些目标又有许多原像，立即说明唯一解需要额外条件。一般公式中的非零分母常常对应这种几何退化；检查它，比在算完后才发现除以零更有效。



\paragraph{3. 用不变量与简单例子回查答案}

切到剪切，形状改变但面积倍率为一。若计算结果声称面积翻倍，就可以迅速定位错误。解题结束不只是写出一个表达式，还应代回原条件，核对单位、符号和特殊值，并说明结论在哪些情况下失效。图形提供检验线索，证明负责说明一般性。
